\subsection{Từ bàn cờ sang ràng buộc và biến CNF}

Trong mã nguồn, biến ô $(r,c)$ được đánh số
\begin{equation}
  v(r,c)=rN+c+1,\qquad 0\leq r,c<N.
\end{equation}
Đánh số từ 1 phù hợp với quy ước literal của PySAT; literal âm biểu thị phủ
định. Với $N=4$, nghiệm có các ô
$(0,1),(1,3),(2,0),(3,2)$ tương ứng các biến $2,8,9,15$. Đây là một ví
dụ kiểm tra chuyển đổi chỉ số, không phải nghiệm duy nhất của $N=4$.

Mỗi hàng và mỗi cột nhận một clause ALO. Chẳng hạn, hàng $r$ có
\begin{equation}
  (x_{r,0}\lor x_{r,1}\lor\cdots\lor x_{r,N-1}).
\end{equation}
AMO trên cùng nhóm ngăn hai ô cùng đúng; ALO cộng AMO tạo exactly-one.
Với đường chéo, chỉ cần AMO. Hai họ đường chéo được nhận biết bởi hằng số
$r-c$ và $r+c$. Trong cài đặt SAT, đường chéo dài một ô được bỏ qua vì AMO
trên một literal luôn đúng; MIP vẫn tạo bất đẳng thức tương ứng nhưng nó dư.

Với $N\geq2$, có $2N$ nhóm hàng/cột và $4N-6$ đường chéo dài ít nhất
hai ô cần AMO. Nếu dùng Pairwise cho mọi nhóm, số clause AMO có thể tính
trực tiếp:
\begin{equation}
 C_{\mathrm{pair}}(N)
 =2N\binom{N}{2}+4\binom{N}{3}+2\binom{N}{2}.
\end{equation}
Hạng đầu thuộc hàng/cột. Mỗi hướng chéo có dãy độ dài
$1,2,\ldots,N-1,N,N-1,\ldots,1$; cộng số cặp trên cả hai hướng cho hai hạng
còn lại. Sau khi thêm $2N$ clause ALO, $N=8$ cho
$C_{\mathrm{pair}}(8)+16=744$ clause, đúng bằng số do chương trình ghi nhận.
Đây là phép kiểm tra chéo giữa công thức và implementation, đồng thời giải
thích vì sao Pairwise tăng bậc ba theo $N$.

\subsection{Đồng nhất nghiệm giữa ba biểu diễn}

Đặt $q_r=c$ khi và chỉ khi $x_{r,c}=1$. Exactly-one trên từng hàng làm ánh
xạ này xác định được. Exactly-one trên cột tương đương các $q_r$ đôi một
khác nhau. Hai quân ở hàng $r\neq s$ nằm chung đường chéo chính khi
$q_r-r=q_s-s$, và chung đường chéo phụ khi $q_r+r=q_s+s$.
Vì vậy ba ràng buộc \texttt{AllDifferent} của CP, các nhóm AMO/ALO của
SAT và các phương trình/bất phương trình MIP diễn tả cùng tập bàn cờ.

Hàm \texttt{validate\_solution} không cần tin vào trạng thái do solver báo.
Nó đòi đúng $N$ tọa độ trong miền, rồi kiểm tra bốn tập khóa: hàng, cột,
$r-c$ và $r+c$. Cách kiểm tra độc lập này phát hiện cả nghiệm trùng ô,
sai miền và xung đột đường chéo. Nó được gọi sau mọi lần solver báo SAT;
run sai sẽ có trạng thái \texttt{INVALID\_SOLUTION} thay vì được tính là
thành công.
